%; whizzy paragraph -pdf xpdf -latex ./whizzypdfptex.sh
%; whizzy-paragraph "^\\\\begin{frame}"
% latex beamer presentation.
% platex, latex-beamer でコンパイルすることを想定。 

%     Tokyo Debian Meeting resources
%     Copyright (C) 2009 Junichi Uekawa
%     Copyright (C) 2009 Nobuhiro Iwamatsu

%     This program is free software; you can redistribute it and/or modify
%     it under the terms of the GNU General Public License as published by
%     the Free Software Foundation; either version 2 of the License, or
%     (at your option) any later version.

%     This program is distributed in the hope that it will be useful,
%     but WITHOUT ANY WARRANTY; without even the implied warreanty of
%     MERCHANTABILITY or FITNESS FOR A PARTICULAR PURPOSE.  See the
%     GNU General Public License for more details.

%     You should have received a copy of the GNU General Public License
%     along with this program; if not, write to the Free Software
%     Foundation, Inc., 51 Franklin St, Fifth Floor, Boston, MA  02110-1301 USA

\documentclass[cjk,dvipdfmx,12pt]{beamer}
\usetheme{Tokyo}
\usepackage{monthlypresentation}

%  preview (shell-command (concat "evince " (replace-regexp-in-string
%  "tex$" "pdf"(buffer-file-name)) "&")) 
%  presentation (shell-command (concat "xpdf -fullscreen " (replace-regexp-in-string "tex$" "pdf"(buffer-file-name)) "&"))
%  presentation (shell-command (concat "evince " (replace-regexp-in-string "tex$" "pdf"(buffer-file-name)) "&"))

%http://www.naney.org/diki/dk/hyperref.html
%日本語EUC系環境の時
\AtBeginDvi{\special{pdf:tounicode EUC-UCS2}}
%シフトJIS系環境の時
%\AtBeginDvi{\special{pdf:tounicode 90ms-RKSJ-UCS2}}

\newenvironment{commandlinesmall}%
{\VerbatimEnvironment
  \begin{Sbox}\begin{minipage}{1.0\hsize}\begin{fontsize}{8}{8} \begin{BVerbatim}}%
{\end{BVerbatim}\end{fontsize}\end{minipage}\end{Sbox}
  \setlength{\fboxsep}{8pt}
% start on a new paragraph

\vspace{6pt}% skip before
\fcolorbox{dancerdarkblue}{dancerlightblue}{\TheSbox}

\vspace{6pt}% skip after
}
%end of commandlinesmall

\title{第168回東京エリアDebian勉強会 \\　\\debianにおけるnginxの設定例}
\subtitle{}
\author{Norimitsu Sugimoto (杉本 典充) \\dictoss@live.jp}
\date{2018-11-17}
\logo{\includegraphics[width=8cm]{image200607/openlogo-light.eps}}

\begin{document}

\section{表紙}

\begin{frame}
  \titlepage{}
\end{frame}


\section{自己紹介}

\begin{frame}{自己紹介}
  \begin{itemize}
  \item Norimitsu Sugimoto (杉本 典充)
  \item dictoss@live.jp
  \item Twitter: @dictoss
  \item Debian使って15年以上。sargeがtestingの頃から使っている
  \item 仕事はソフトウェア開発者をやってます
  \item pythonとDjangoの組み合わせで使うことが多いです
  \end{itemize}
\end{frame}


\section{目次}

\begin{frame}{アジェンダ}
  \begin{itemize}
  \item はじめに
  \item debianにおけるwebサーバのパッケージ
  \item nginx
  \item nginxの設定例
  \item まとめ・参考文献
  \end{itemize}
\end{frame}


\section{debianにおけるwebサーバのパッケージ}
\emtext{debianにおけるwebサーバのパッケージ}

\begin{frame}[containsverbatim]{debianにおけるwebサーバのパッケージ}
  \begin{itemize}
  \item webサーバの利用調査
    \begin{itemize}
    \item \url{https://w3techs.com/technologies/comparison/ws-apache,ws-microsoftiis,ws-nginx}
    \end{itemize}
  \item debianではapache、nginxのパッケージを用意されており、aptコマンドでインストール可能
  \end{itemize}
\end{frame}


\section{nginx}
\emtext{nginx}

\begin{frame}[containsverbatim]{nginxとは}
  \begin{itemize}
  \item webサイト \url{https://nginx.org/}
  \item エンジンエックスと読みます
  \item ソースコードは、2-clause BSD-like licenseで公開
  \item 最初のリリースは2004年
  \item イベント駆動のシステムコールを利用（linuxではepollを利用）し、大量のクライアント（C10K）を捌く
  \item HTTP/HTTPSのファイル配信のみならず、HTTP/HTTPSリバースプロキシ、mailプロキシ、TCP/UDPプロキシとして動作
  \end{itemize}
\end{frame}

\begin{frame}[containsverbatim]{debianにおけるnginxのインストール方法とディレクトリ構成}
aptでインストール可能
\begin{commandline}
# apt install nginx
\end{commandline}

拡張モジュールの数でdebパッケージを3つに分かれている
\begin{itemize}
  \item nginx-full (=nginx)
  \item nginx-light
  \item nginx-extras
\end{itemize}

\end{frame}


\section{nginxの設定例}
\emtext{nginxの設定例}

\subsection{nginxパッケージのインストール直後の設定}
\begin{frame}[containsverbatim]{nginxパッケージのインストール直後の設定}

一部省略して抜粋 
\begin{commandline}
$ tree /etc/nginx
  |--  conf.d/
  |--  fastcgi_params
  |--  modules-available/
  |--  modules-enable/
  |--  nginx.conf
  |--  proxy_params
  |--  site-available/
  |--  site-enabled/
  |--  snippets/
  |--  wsgi_params
\end{commandline}
\end{frame}

\begin{frame}[containsverbatim]{nginxパッケージのインストール直後の設定}

一部省略して抜粋
\begin{commandline}
  $ cat /etc/nginx/nginx.conf
  user www-data;
  worker_processes auto;
  include /etc/nginx/modules-enabled/*.conf;
  events {
    worker_connections 768;
  }
  http {
    access_log /var/log/nginx/access.log;
    error_log /var/log/nginx/error.log;
    gzip on;
    gzip_disable ``msie6'';
    include /etc/nginx/conf.d/*.conf;
    include /etc/nginx/sites-enabled/*;
  }
\end{commandline}
\end{frame}

\begin{frame}[containsverbatim]{VirtualHost設定"default"}
  \begin{commandline}
$ cat /etc/nginx/sites-available/default
server {
  listen 80 default_server;
  listen [::]:80 default_server;
  root /var/www/html;
  index index.html index.htm index.nginx-debian.html;
  server_name _;
  location / {
    try_files $uri $uri/ =404;
  }
}
\end{commandline}
\end{frame}


\subsection{SSL/TLS設定}
\emtext{SSL/TLS設定}

\begin{frame}[containsverbatim]{SSL/TLS設定}
  \begin{itemize}
  \item /etc/nginx/sites-available/ssl にSSL/TLS待ち受け設定を書く
  \item /etc/nginx/sites-enabled/ssl にシンボリックリンクを生成し、VirtualHostを追加
  \item 【参考情報】Mozillaが公開している設定の生成ツール
    \begin{itemize}
      \item \url{https://mozilla.github.io/server-side-tls/ssl-config-generator/}
    \end{itemize}
  \end{itemize}
\end{frame}

\begin{frame}[containsverbatim]{SSL/TLS設定}
設定のポイント
\begin{itemize}
\item 中間証明書を必要とする場合は１つのcrtファイルに結合する
\item ssl\_protocols TLSv1.2; とだけ書くと、TLSv1.2のみで接続を受け付ける
\item \url{https://www.ssllabs.com/ssltest/} などでスキャンして診断
\item ssl\_dhparamは DH 1024bitより上げるか悩みどころ
\end{itemize}
\end{frame}

\begin{frame}[containsverbatim]{SSL/TLS設定}

一部抜粋
\begin{commandlinesmall}
$ cat /etc/nginx/sites-available/ssl
  server {
    listen 443 ssl http2;
    listen [::]:443 ssl http2;
    ssl on;
    ssl_certificate /etc/ssl/private/server.crt;
    ssl_certificate_key  /etc/ssl/private/server.key;
    ssl_protocols TLSv1.2; # Dropping SSLv3, ref: POODLE
    ssl_prefer_server_ciphers on; # ref: FREAK
    ssl_ciphers  HIGH:!aNULL:!MD5;
    gzip off;  # unuse gzip, ref: BREACH
    root /var/www/html;
    server_name www.example.com;
    location / {
      try_files $uri $uri/ =404;
    }
  }
\end{commandlinesmall}
\end{frame}

\begin{frame}[containsverbatim]{SSL/TLS設定：適用}
\begin{commandlinesmall}
$ sudo nginx -t
nginx: the configuration file /etc/nginx/nginx.conf syntax is ok
nginx: configuration file /etc/nginx/nginx.conf test is successful

$ sudo systemctl restart nginx
\end{commandlinesmall}

設定したサーバの"https://hostname/" へアクセスする
\end{frame}


\subsection{リバースproxy設定}
\emtext{リバースproxy設定}

\begin{frame}[containsverbatim]{リバースproxy}
  \begin{itemize}
  \item リバースproxy
    \begin{itemize}
    \item クライアントと直接接続するのはnginxが担う
    \item クライアントへ応答する実際のデータ生成はバックエンドのサーバが行うように通信を中継する機能
    \end{itemize}
  \item SSL/TLS設定をするとき、リバースproxyのnginxのみにSSL証明書を配置すればよい。バックエンドのサーバはHTTP設定のみでよい。
  \end{itemize}
\end{frame}

\begin{frame}[containsverbatim]{リバースproxy：upstreamの設定}
  \begin{itemize}
  \item upstream句にバックエンドのserverの通信先を設定する
  \item 複数行のserver設定をすると、振り分けできる
  \item 重み付きラウンドロビン、leas\_conn、ip\_hashが利用可能
  \end{itemize}

\begin{commandline}
$ cat /etc/nginx/conf.d/upstream_proxy.conf
  upstream backend_app1 {
    # least_conn;
    # ip_hash;

    server 192.168.1.100:80 weight=1;
    server 192.168.1.101:80 weight=1;
  }    
\end{commandline}
\end{frame}

\begin{frame}[containsverbatim]{リバースproxy：VirtualHost設定}
  \begin{itemize}
  \item proxy先への通信はデフォルトでHTTP/1.0であり、"proxy\_http\_version"で変更可能
  \end{itemize}

\begin{commandlinesmall}
$ cat /etc/nginx/sites-available/default
server {
  # (snip)
  location ~ ^/proxy/(.*)$ {
    proxy_pass http://backend_proxy1;
    
    # proxy_http_version 1.1;

    proxy_set_header X-Real-IP $remote_addr;
    proxy_set_header X-Forwarded-Proto $scheme;
    proxy_set_header X-Forwarded-For $proxy_add_x_forwarded_for;
    proxy_set_header Host $http_host;
    #proxy_redirect http://backend_proxy1/ http://www.example.com/;
  }
  # (snip)
\end{commandlinesmall}
\end{frame}


\begin{frame}[containsverbatim]{リバースproxy：適用}
\begin{commandlinesmall}
$ sudo nginx -t
nginx: the configuration file /etc/nginx/nginx.conf syntax is ok
nginx: configuration file /etc/nginx/nginx.conf test is successful

$ sudo systemctl restart nginx
\end{commandlinesmall}

設定したサーバの"/proxy/" へアクセスする
\end{frame}


\subsection{FastCGI設定}
\emtext{FastCGI設定}

\begin{frame}[containsverbatim]{FastCGI：PHP-FPMのインストール}

PHP-7.0とPHP-7.0用PHP-FPMをインストール
\begin{commandline}
$ sudo apt install php7.0 php7.0-fpm
\end{commandline}

待ち受けをinet socketに変更
\begin{commandline}
$ sudo vi /etc/php/7.0/fpm/pool.d/www.conf
;listen = /run/php/php7.0-fpm.sock
listen = 9000
\end{commandline}

PHP-FPMを再起動
\begin{commandline}
$ sudo systemctl restart php7.0-fpm
$ ss -npta | grep 9000
LISTEN     0  128  :::9000   :::*
\end{commandline}
\end{frame}

\begin{frame}[containsverbatim]{FastCGI：PHPスクリプト配置}
phpinfo()を実行するスクリプトを動作確認用に配置
\begin{commandline}
$ sudo mkdir /var/www/html/myphpapp
$ sudo vi /var/www/html/myphpapp/phpinfo.php
<?php
    phpinfo();
\end{commandline}
\end{frame}

\begin{frame}[containsverbatim]{FastCGI：upstream設定}
\begin{commandline}
$ sudo vi /etc/nginx/conf.d/upstream_fcgi.conf
upstream backend_fcgi1 {
  # least_conn;
  # ip_hash;

  server 127.0.0.1:9000;
}
\end{commandline}
\end{frame}

\begin{frame}[containsverbatim]{FastCGI：VirtualHost設定}

".php"へのアクセスをすべて、FastCGIサーバへ中継する
\begin{commandline}
$ sudo vi /etc/nginx/sites-available/default
  # (snip)
  location ~ \.php$ {
    include snippets/fastcgi-php.conf;

    fastcgi_pass backend_fcgi1;
  }
  # (snip)
\end{commandline}
\end{frame}

\begin{frame}[containsverbatim]{FastCGI：適用}
\begin{commandlinesmall}
$ sudo nginx -t
nginx: the configuration file /etc/nginx/nginx.conf syntax is ok
nginx: configuration file /etc/nginx/nginx.conf test is successful

$ sudo systemctl restart nginx
\end{commandlinesmall}

設定したサーバの"/myphpapp/phpinfo.php" へアクセスする
\end{frame}


\subsection{WSGI設定}
\emtext{WSGI設定}

\begin{frame}[containsverbatim]{WSGI：uwsgiのインストール}

python3.5、pip3、uwsgiをインストール
\begin{commandline}
$ sudo apt install python3 python3-pip \
 uwsgi uwsgi-plugin-python3
\end{commandline}

\end{frame}

\begin{frame}[containsverbatim]{WSGI：アプリケーション配備}

\url{https://github.com/dictoss/django-tutorial} を実験用に配置

\begin{commandlinesmall}
$ sudo apt install git
$ sudo pip3 install -U django==2.0.9

$ cd
$ git clone https://github.com/dictoss/django-tutorial.git
$ sudo mkdir /var/www/wsgi_apps_uwsgi
$ sudo cp -r django-tutorial/2.0/mysite /var/www/wsgi_apps_uwsgi/
$ sudo chown -fR www-data:www-data /var/www/wsgi_apps_uwsgi/mysite
$ ls /var/www/wsgi_apps_uwsgi/mysite
db.sqlite3  manage.py  mysite  polls
\end{commandlinesmall}
\end{frame}

\begin{frame}[containsverbatim]{WSGI：uwsgi設定}

設定を一部抜粋
\begin{commandline}
$ sudo vi /etc/uwsgi/apps-available/django-tutorial.ini
[uwsgi]
uid = www-data
gid = www-data
plugin-dir = /usr/lib/uwsgi/plugins
plugin = python3
base = /var/www/wsgi_apps_uwsgi/mysite
chdir = /var/www/wsgi_apps_uwsgi/mysite
module = mysite.wsgi
callable = application
env =
socket = 0.0.0.0:3031
processes = 2
threads = 32
master = True
vacuum = True
harakiri = 60
max-requests = 512
\end{commandline}
\end{frame}

\begin{frame}[containsverbatim]{WSGI：uwsgi設定}

apps-availableのiniファイルをapps-enabledへシンボリックリンクを生成する
\begin{commandline}
$ cd /etc/uwsgi/apps-enabled
$ sudo ln -fs ../apps-available/django-tutorial.ini .
$ tree /etc/uwsgi
  /etc/uwsgi
  ├──  apps-available
  │ ├──  README
  │ ├──  django-tutorial.ini
  └──  apps-enabled
  ├──  README
  └──  django-tutorial.ini ->
              ../apps-available/django-tutorial.ini
\end{commandline}
\end{frame}

\begin{frame}[containsverbatim]{WSGI：uwsgi設定}

uwsgiを再起動して設定反映
\begin{commandline}
$ sudo systemctl restart uwsgi
$ ss -npta | grep 3031
LISTEN     0  100  *:3031  *:*
\end{commandline}
\end{frame}

\begin{frame}[containsverbatim]{WSGI：upstream設定}
\begin{commandline}
$ sudo vi /etc/nginx/conf.d/upstream_uwsgi.conf
upstream backend_uwsgi1 {
  # least_conn;
  # ip_hash;

  server 127.0.0.1:3031;
}
\end{commandline}
\end{frame}

\begin{frame}[containsverbatim]{WSGI：VirtualHost設定}

"/mysite/"配下へのアクセスをすべて、WSGIサーバへ中継する
\begin{commandline}
$ sudo vi /etc/nginx/sites-available/default
  # (snip)
  location ~ ^/mysite/(.*)$ {
    include uwsgi_params;

    uwsgi_param SCRIPT_NAME /mysite;
    uwsgi_param PATH_INFO /$1;

    uwsgi_pass backend_uwsgi1;
  }
  # (snip)
\end{commandline}
\end{frame}

\begin{frame}[containsverbatim]{WSGI：適用}
\begin{commandlinesmall}
$ sudo nginx -t
nginx: the configuration file /etc/nginx/nginx.conf syntax is ok
nginx: configuration file /etc/nginx/nginx.conf test is successful

$ sudo systemctl restart nginx
\end{commandlinesmall}

設定したサーバの"/mysite/polls/" へアクセスする
\end{frame}


\section{まとめ・参考文献}
\emtext{まとめ・参考文献}

\begin{frame}[containsverbatim]{まとめ・参考文献}

まとめ
\begin{itemize}
  \item nginxのよく使う設定をまとめてみました
  \item 大規模向けに性能向上するパラメータについてはチューニングしてみてください
  \end{itemize}

  参考文献
  \begin{itemize}
  \item nginx documentation \url{https://nginx.org/en/docs/}
  \item nginx - DebianWiki \url{https://wiki.debian.org/Nginx}
  \end{itemize}
\end{frame}

\end{document}


;;; Local Variables: ***
;;; outline-regexp: "\\([ 	]*\\\\\\(documentstyle\\|documentclass\\|emtext\\|section\\|begin{frame}\\)\\*?[ 	]*[[{]\\|[]+\\)" ***
;;; End: ***
